\documentclass[12pt,a4paper]{article}

\usepackage[utf8]{inputenc}
\usepackage[T1]{fontenc}
\usepackage[brazil]{babel}
\usepackage[margin=2.5cm]{geometry}
\usepackage{graphicx}

\title{$titulo}
\author{$autor \\ Nome da instituição}
\date{}

\begin{document}

\maketitle

\section{Introdução}

Em poucos parágrafos: o tema, o problema e por que ele é importante. Resumos expandidos costumam ter de 2 a 4 páginas, então seja direto.

\section{Objetivos}

O que o trabalho pretende alcançar. Uma frase para o objetivo geral e, se quiser, uma lista curta de objetivos específicos.

\section{Metodologia}

Como o trabalho foi ou será feito.

\section{Resultados}

O que já foi obtido. Se o projeto ainda está em andamento, escreva os resultados esperados.

\section{Conclusão}

O que os resultados mostram e os próximos passos.

\section*{Agradecimentos}

Agradeça a orientadores, colegas e a quem financiou o projeto (bolsa, edital).

\bibliographystyle{plain}
\nocite{barbosa2010ihc}
\bibliography{referencias}

\end{document}
